% !TeX spellcheck = ru_RU
% !TEX root = vkr.tex

\section{Введение}

Одним из центральных направлений развития систем на основе больших языковых моделей (LLM) является создание автономных агентов, способных взаимодействовать с внешними инструментами~--- API, базами данных, поисковыми системами~--- для решения задач, выходящих за пределы генерации текста. Подход, при котором модель самостоятельно определяет, какой инструмент вызвать и с какими параметрами, получил название \textit{tool use} (или \textit{function calling}) и активно исследуется в работах~\cite{schick2023toolformer, patil2023gorilla, qin2024toollearning}. Качество выбора инструмента непосредственно определяет применимость агента в промышленных системах: ошибка классификации приводит либо к некорректному ответу, либо к полному отсутствию результата.

В настоящей работе рассматривается AI-агент в сфере недвижимости, обрабатывающий пользовательские запросы на естественном языке. Агент должен классифицировать каждый запрос в один из двух классов и вызвать соответствующую функцию: \texttt{ResolveEntities} для географических объектов (класс GEO) и \texttt{SearchTags} для характеристик объектов недвижимости (класс TAG). Базовая модель YandexGPT~5.1~Pro~\cite{yandexgpt2024docs}, используемая в качестве ядра агента, на объединённом бенчмарке из 500 примеров достигает точности (accuracy) лишь $81{,}8\%$ при Macro~$F_1 = 0{,}526$. Анализ ошибок выявляет систематический сдвиг (bias): из 174 запросов класса TAG верно классифицированы лишь $49{,}4\%$, а для 34 запросов модель вообще не генерирует вызов инструмента. Таким образом, базовая модель обладает выраженным предпочтением класса GEO, что делает её непригодной для промышленного применения без дополнительной адаптации.

Центральная рабочая гипотеза данного исследования состоит в том, что выбор инструмента может быть интерпретирован как устойчивый \textit{паттерн формы ответа} модели~--- по аналогии с гипотезой поверхностного выравнивания (Superficial Alignment Hypothesis) Zhou и~др.~\cite{zhou2023lima, gudibande2023false}, согласно которой дообучение с инструкциями преимущественно корректирует формат и стиль генерации, а не фактические знания модели (подробно~--- раздел~\ref{sec:related}). Настоящая работа проверяет применимость этой аналогии к задаче tool selection: если она верна, то даже небольшой объём обучающих примеров, дообученных методом LoRA~\cite{hu2022lora}, должен радикально улучшить качество классификации.

Целью работы является экспериментальная проверка указанной гипотезы. Для этого решаются следующие задачи:
\begin{enumerate}
    \item формализация задачи классификации инструментов как задачи бинарной классификации с двумя бенчмарками (GEO и TAG);
    \item разработка инфраструктуры для дообучения и систематической оценки моделей;
    \item оценка 12 конфигураций моделей (1 базовая и 11 дообученных) с вариацией типа обучающих данных (GEO, TAG, смешанные) и объёма выборки;
    \item статистический анализ результатов с применением критериев Кохрена~$Q$ и Макнемара для попарных сравнений моделей.
\end{enumerate}

Научная новизна работы заключается в исследовании применения метода LoRA для коммерческой LLM на специфичном русскоязычном гео-домене с использованием строгих статистических критериев оценки. В отличие от существующих работ, сосредоточенных преимущественно на англоязычных моделях открытого доступа, данное исследование включает строгое статистическое сравнение 12~конфигураций моделей на двух независимых бенчмарках. Основной результат~--- достижение точности $99{,}4\%$ лучшей дообученной моделью при всего одной ошибке на 500 примерах~--- согласуется с рабочей гипотезой о стилистической природе выбора инструмента и демонстрирует практическую эффективность подхода.

Работа организована следующим образом. В разделе~2 дана формальная постановка задачи. Раздел~3 содержит обзор литературы по tool use, методам дообучения и гипотезе поверхностного выравнивания. В разделе~4 описаны теоретические основы используемых методов. Раздел~5 посвящён методологии экспериментов, раздел~6~--- анализу результатов. Раздел~7 обсуждает угрозы валидности полученных результатов. В разделе~8 описана программная реализация. Раздел~9 содержит заключение и направления дальнейших исследований.
